% Template: erratum to a paper published in a journal, drafted from the errors
% that a Lean formalization found. Copy it to the project, fill it in, compile
% it with `latexmk -pdf ERRATUM.tex`, and present the PDF to the user as a
% draft; the skills never send it. Locations refer to the published version,
% not to the arXiv version that was formalized. List only errors in the paper, as classified in the
% formalization report, each audited independently; keep every item short.
% At the user's request, the same template serves for an erratum to a cited
% work of which the user is an author, from the confirmed entries of
% CITED_WORKS_ERRATA.md.
\documentclass[11pt]{amsart}
\usepackage[T1]{fontenc}
\usepackage{amssymb}
\usepackage[hidelinks]{hyperref}

\title[Erratum]{Erratum to ``[Title of the paper]''}
\author{[Authors of the paper]}
\date{}

\begin{document}

\begin{abstract}
A formal verification in the Lean proof assistant of the results of
[reference: authors, title, journal, volume (year), pages,
\href{https://doi.org/[DOI]}{doi:[DOI]}] found the errors listed below.
[None of them affects the main results. / The main results hold as corrected
below.] The formalization is available at [repository and revision].
\end{abstract}

\maketitle

% One item per error: location in the published version, the correction, then
% the reason. Write the correction as the paper should read.
\begin{enumerate}
\item \emph{[Theorem 2.3], [p.~12, l.~5].}
  [The statement requires the additional hypothesis $\mu(A) > 0$:]
  [the corrected statement or step.]
  [Reason: a counterexample, or the step that fails and why.]
\end{enumerate}

\end{document}
